\section{Exact original coordinates on the actual unit-ideal Rees chart}

\begin{definition}[The actual unit Rees chart retains the original ring]
\label{def:fc-original-unit-rees-chart-coefficients}
\leanok
\lean{EnvelopeProperty.UnitIdealReesCoordinates.fractionCoefficient_surjective, EnvelopeProperty.UnitIdealReesCoordinates.fractionCoefficientEquiv, EnvelopeProperty.UnitIdealReesCoordinates.homogeneousCoefficientEquiv, EnvelopeProperty.UnitIdealReesCoordinates.homogeneousCoefficientEquiv_apply}
For the original domain $A$, the actual coefficient map identifies
$A$ with the same genuine homogeneous Rees chart $A[A/1]$ at the
original element $T$. Both the actual fraction subalgebra and the
actual homogeneous localization are identified through their
original coefficient maps, without a supplied comparison.
\end{definition}
\begin{proof}
The genuine original fraction-chart denominator-clearing theorem at
$g=1$ says that every original chart element is the image of an
actual element of $A$. The actual coefficient map is injective in
the original fraction field, so is an algebra isomorphism. Compose
with the genuine original homogeneous-localization to fraction-chart
isomorphism.
\end{proof}

\begin{theorem}[The same original Rees coordinate ratio]
\label{thm:fc-original-unit-rees-coordinate-ratio}
\leanok
\lean{EnvelopeProperty.UnitIdealReesCoordinates.homogeneousCoefficientEquiv_symm_generator}
\uses{def:fc-original-unit-rees-chart-coefficients}
Under this actual original coefficient identification, the genuine
original homogeneous ratio $(aT)/T$ is exactly the SAME original
$a\in A$ for every $a\in A$.
\end{theorem}
\begin{proof}
Map the actual homogeneous fraction into the same original fraction
field by the genuine Rees chart evaluation. Its value is exactly
$a/1=a$. Both coefficient identifications preserve the literal
original coefficient, and their injectivity gives the equality in
the whole original chart ring.
\end{proof}

\begin{theorem}[The actual inverse of the original unit blowup projection]
\label{thm:fc-original-unit-rees-chart-section}
\leanok
\lean{EnvelopeProperty.UnitIdealReesCoordinates.inverse_fractionChart, EnvelopeProperty.UnitIdealReesCoordinates.inverse_homogeneousChart}
\uses{def:fc-original-unit-rees-chart-coefficients}
The inverse of the actual original unit-ideal blowup projection is
the original chart inclusion composed with the spectrum map of the
preceding original coefficient identification. This holds both in
the actual fraction-chart and actual homogeneous-chart presentations.
\end{theorem}
\begin{proof}
The actual chart inclusion followed by the actual original blowup
projection is the spectrum of the original coefficient map. The
preceding true inverse coefficient map therefore makes the composite
the identity on the SAME original spectrum. The original blowup
projection is already a genuine isomorphism, so its actual inverse
is this chart section. Expanding the genuine original fraction-chart
inclusion yields the identical homogeneous-chart formula.
\end{proof}
